\section{The power--log scale and its exact arithmetic}
\label{sec:scale}

\subsection{Blocks and their order}

Throughout, $w$ denotes a positive real variable tending to $0^+$, $L=\log w\to-\infty$ and $\M(w)=1+|\log w|$.

\begin{definition}[Blocks]
A \emph{power--log block} is an expression $w^\beta P(L)$ with $\beta\in\R$ and $P\in\R[L]$, $P\ne0$. Its \emph{exponent} is $\beta$ and its \emph{logarithmic degree} is $\deg P$.
\end{definition}

\begin{lemma}[Power beats logarithm]
\label{lem:dominance}
For every $\varepsilon>0$ and every real $b$, $w^\varepsilon\M(w)^b\to0$ as $w\to0^+$. Consequently, if $\beta<\gamma$ and $P,Q\ne0$, then $w^\gamma Q(L)=\oo\bigl(w^\beta P(L)\bigr)$; and within a fixed exponent, a higher logarithmic degree dominates: $w^\beta L^j=\oo(w^\beta L^k)$ for $j<k$.
\end{lemma}
\begin{proof}
Put $w=e^{-t}$, $t\to+\infty$: $w^\varepsilon\M(w)^b=e^{-\varepsilon t}(1+t)^b\to0$. For the second claim the ratio is $w^{\gamma-\beta}Q(L)/P(L)$, and $|Q(L)/P(L)|\le C\M(w)^{\deg Q}$ for $w$ small because $P(L)$ has no zeros for $L$ sufficiently negative and $|P(L)|\ge c|L|^{\deg P}\ge c$ there; the first claim finishes. The third claim is clear.
\end{proof}

Blocks are therefore totally ordered by size: by increasing exponent, and at equal exponents by decreasing logarithmic degree. A finite sum of blocks with distinct exponents is asymptotic to its block of least exponent (whole polynomial included), and vanishes identically only if every block does. This is why truncations are taken by \emph{complete blocks}: cutting a block inside its polynomial would drop terms larger than the retained terms of the next exponent.

\begin{definition}[Remainder classes]
For $\beta\in\R$ and an integer $k\ge0$ we write $R(w)=\OO_{\beta,k}$, or in the package \wl{PowerLogRemainder[w, $\beta$, k]}, for $R(w)=\OO\bigl(w^\beta\M(w)^k\bigr)$ as $w\to0^+$.
\end{definition}

By Lemma~\ref{lem:dominance}, $\OO_{\beta,k}\subseteq\OO_{\beta',k'}$ whenever $\beta>\beta'$, or $\beta=\beta'$ and $k\le k'$; the class with the smaller exponent, and at equal exponents the larger degree, is the coarser one. A block $w^\beta P(L)$ with $\deg P=k$ belongs to $\OO_{\beta,k}$ and to no $\OO_{\beta,k'}$ with $k'<k$. Sums and products of remainder classes obey
\begin{equation}
 \OO_{\beta,k}+\OO_{\beta',k'}=\OO_{\min(\beta,\beta'),\,k''},\qquad
 \OO_{\beta,k}\cdot\OO_{\beta',k'}=\OO_{\beta+\beta',\,k+k'},
 \label{eq:O-arithmetic}
\end{equation}
where $k''$ is the degree attached to the smaller exponent, and the larger of the two when the exponents coincide.

\subsection{Exponent semigroups}

\begin{lemma}[Local finiteness]
\label{lem:finite}
Let $\delta_1,\dots,\delta_m>0$ and $S=\{\bk\cdot\bdelta:\bk\in\N^m\}$. For every real $H$ the set $\{\bk\in\N^m:\bk\cdot\bdelta\le H\}$ is finite, hence $S\cap(-\infty,H]$ is finite, and every $\sigma\in S$ has only finitely many representations $\sigma=\bk\cdot\bdelta$.
\end{lemma}
\begin{proof}
If $\bk\cdot\bdelta\le H$ then $k_i\le H/\delta_i$ for each $i$.
\end{proof}

When the gaps are incommensurable the \emph{group} generated by them is dense in $\R$ (for instance $\Z+\sqrt2\Z$), but the semigroup of nonnegative combinations has only finitely many elements below any bound. This is the whole reason why expansions with irrational exponents can be computed term by term: the support of every expansion in this article lies in a translate $\beta_0+S$ of such a semigroup. Following \cite{proveit}, a subset of $\R$ that is finite below every bound is called \emph{left-finite}; the semigroups $S$ are left-finite, and so are their translates and finite unions.

\subsection{Jets and their arithmetic}

\begin{definition}[Jets]
Fix gaps as above and $\beta_0\in\R$. A \emph{jet} (of cutoff $H$) is a finite sum $U=\sum_{\sigma\in T}w^\sigma C_\sigma(L)$ with $T\subset\beta_0+S$ finite, $C_\sigma\in\R[L]$, all $\sigma<H$; equal exponents are merged and zero polynomials are discarded. Its \emph{valuation} $\val U$ is its least exponent ($+\infty$ for the empty jet). The jet algebra $\PL_{<H}$ consists of jets truncated below $H$, with the sum and the product computed exactly and then truncated.
\end{definition}

The jet algebra is the quotient of the ring of finite sums of blocks by the ideal of blocks of exponent $\ge H$. In the package a jet is a sorted list of pairs $\{\sigma,C_\sigma\}$ with $\sigma$ an exact real number in canonical form (\wl{RootReduce} for algebraic numbers) and $C_\sigma$ a polynomial in a private symbol $\ell$; two exponents are equal if and only if their canonical forms coincide, and they are ordered exactly (Section~\ref{sec:package}).

\subsubsection*{Differentiation}

\begin{lemma}[Euler identity]
\label{lem:euler}
For $\beta\in\R$, $P\in\R[L]$ and $r\in\N$,
\begin{equation}
 \frac{\dd}{\dd w}\bigl[w^\beta P(L)\bigr]=w^{\beta-1}(\beta+\D)P(L),\qquad
 \frac{\dd^r}{\dd w^r}\bigl[w^\beta P(L)\bigr]=w^{\beta-r}\prod_{j=0}^{r-1}(\beta-j+\D)P(L).
 \label{eq:euler}
\end{equation}
The operators $\beta-j+\D$ commute, and $(\beta+\D)$ preserves $\deg P$ when $\beta\ne0$ and lowers it by one when $\beta=0$.
\end{lemma}
\begin{proof}
$\frac{\dd}{\dd w}[w^\beta P(\log w)]=\beta w^{\beta-1}P+w^\beta P'(\log w)/w$. Iterate.
\end{proof}

The identity says that differentiation acts on a block as the scalar $\beta$ plus $\D$ on the polynomial; the exponent bookkeeping and the polynomial arithmetic separate completely. In the coefficient formulas below the operators $\D+c$ with $c>0$ appear, and their action on the coefficient vector $(z_0,\dots,z_d)$ of $P=\sum z_sL^s$ is bidiagonal: $(c+\D)P=\sum_s\bigl(cz_s+(s+1)z_{s+1}\bigr)L^s$.

\subsubsection*{Unit series}

Let $U$ be a jet with $\val U>0$. Then $U^k$ has valuation $\ge k\val U$, so in $\PL_{<H}$ only finitely many powers of $U$ are nonzero, and every power series $\Phi(t)=\sum_k\phi_kt^k$ can be substituted:
\begin{equation}
 \Phi(U)=\sum_{k<H/\val U}\phi_kU^k\in\PL_{<H}.
 \label{eq:unit-series}
\end{equation}
In particular $(1+U)^a=\sum_k\binom akU^k$ for every real $a$ (with $\binom ak=a(a-1)\cdots(a-k+1)/k!$), $\log(1+U)=\sum_{k\ge1}(-1)^{k+1}U^k/k$, $e^U=\sum_kU^k/k!$, and $\phi(c+U)=\sum_k\phi^{(k)}(c)U^k/k!$ for any $\phi$ analytic at $c$. These are identities between finite sums of blocks modulo exponents $\ge H$; their analytic content is discussed in Section~\ref{sec:convergence}.

\subsubsection*{Composition with the logarithm}

The one operation that mixes the exponent bookkeeping with the polynomial coefficients is the substitution of a jet into the logarithm: if $w$ is replaced by $w(1+U)$ then $L$ becomes $L+\log(1+U)$, and a block $w^a P(L)$ becomes $w^a(1+U)^aP(L+\log(1+U))$. The following recurrence computes this without expanding $\log(1+U)$ separately.

\begin{lemma}[Composition recurrence]
\label{lem:composition}
Let $P\in\R[L]$ and $a\in\R$. For $|t|<1$,
\begin{equation}
 (1+t)^aP\bigl(L+\log(1+t)\bigr)=\sum_{k\ge0}Q_k(L)\,t^k,\qquad
 Q_0=P,\qquad Q_{k+1}=\frac{(a-k)Q_k+Q_k'}{k+1},
 \label{eq:composition}
\end{equation}
with $Q_k\in\R[L]$, $\deg Q_k\le\deg P$, and, when $a\in\N$ and $P$ is constant, $Q_k=0$ for $k>a$. Consequently, for a jet $U$ of positive valuation, $(1+U)^aP(L+\log(1+U))=\sum_kQ_k(L)U^k$ in $\PL_{<H}$.
\end{lemma}
\begin{proof}
Put $F(t)=(1+t)^aP(L+\log(1+t))$; it is analytic in $|t|<1$ with values in $\R[L]$, so $F=\sum Q_kt^k$ with $Q_k\in\R[L]$ and $Q_0=P$. Differentiating, $(1+t)\partial_tF=aF+\partial_LF$, because $\partial_t$ acting on the two factors gives $a(1+t)^{a-1}P+(1+t)^{a-1}P'$. Comparing coefficients of $t^k$ in $(1+t)\sum(k+1)Q_{k+1}t^k=\sum(aQ_k+Q_k')t^k$ gives $(k+1)Q_{k+1}+kQ_k=aQ_k+Q_k'$. The degree does not increase, and for $a\in\N$ and constant $P$ the factor $(a-k)$ annihilates $Q_{a+1}$ and hence all later $Q_k$ (for nonconstant $P$ the derivative keeps the sequence alive: $\log(1+t)$ is not a polynomial). The jet statement follows because substitution of $U$ for $t$ is a ring homomorphism on finite sums.
\end{proof}

\begin{remark}
The recurrence is the whole composition calculus needed here: the forward model $ax^p\bigl(1+\sum x^{\delta_i}B_i(\log x)\bigr)$ evaluated at $x=z(1+U)$ is
$az^p(1+U)^p\bigl(1+\sum_iz^{\delta_i}(1+U)^{\delta_i}B_i(L+\log(1+U))\bigr)$, each factor of which is an instance of \eqref{eq:composition} (with $P=1$ for the first). Its derivative is obtained from the same lemma with $a$ replaced by $a-1$ and $P$ by $aP+P'$, since $\frac{\dd}{\dd x}[x^aP(\log x)]=x^{a-1}(aP+P')(\log x)$.
\end{remark}

\subsection{Blocks with nonconstant leading coefficient}

Everything above allows the polynomial in the leading block to be nonconstant. Two operations do not: taking a real non-integer power or a logarithm of a sum whose leading block is $w^\beta P(L)$ with $\deg P\ge1$ produces $P(L)^a$ or $\log P(L)$, which are not polynomials in $L$. Such expressions live in larger scales (rational functions of $L$, or iterated logarithms, Section~\ref{sec:boundaries}) and are outside the class treated here; the package reports them as unsupported. Nonnegative integer powers of any jet are computed by repeated multiplication and stay in the class.
